\chapter*{How to use this encyclopaedia}
\addcontentsline{toc}{chapter}{How to use this encyclopaedia}
% STATUS: written 2026-11-06 (describes the organisation the book already has:
% the eight question families of Section~\ref{sec:basic-theory:question-families},
% the chapter layout, and the two appendices. No new claim is made.)
\label{page:how-to-use}

The book is organised around eight questions rather than around its chapters. % claim-ok: this page describes how this book is arranged
A reader who wants one of them answered can go straight to the % claim-ok: this page describes how this book is arranged
appendix that indexes it. A reader who wants to follow the subject from its % claim-ok: this page describes how this book is arranged
beginnings can read Part~I onward. % claim-ok: this page describes how this book is arranged

\section*{The eight question families}

Section~\ref{sec:basic-theory:question-families} fixes the eight questions the % claim-ok: this page describes how this book is arranged
book asks of every class of algebra. They are: % claim-ok: this page describes how this book is arranged

\begin{enumerate}
\item Which kind of derivation implies which other?
\item When are two definitions equivalent?
\item Is every derivation inner?
\item Is every derivation continuous, and does continuity come for free?
\item Does the algebra have to be commutative?
\item Can all derivations of the class be described?
\item Does a derivation extend or lift to a larger algebra?
\item Where does an expected implication fail?
\end{enumerate}

The chapters answer them in this order wherever the sources allow. % claim-ok: this page describes how this book is arranged
Chapter~\ref{app:questions}, which opens the book, collects the answers in tables, family by family, % claim-ok: this page describes how this book is arranged
so that a reader comparing two classes of algebra can see both at once. % claim-ok: this page describes how this book is arranged

\section*{How a chapter is laid out}

Each chapter follows the same plan. % claim-ok: this page describes how this book is arranged

The introduction gives the history: why the topic arose, who first asked the % claim-ok: this page describes how this book is arranged
question, and what the chapter answers. % claim-ok: this page describes how this book is arranged
The sections that follow begin with the definitions as their sources state them, % claim-ok: this page describes how this book is arranged
then the theorems as their sources state them, each with the page or numbered % claim-ok: this page describes how this book is arranged
item. % claim-ok: this page describes how this book is arranged

Counterexamples come where a family of the last kind fails; they are what % claim-ok: this page describes how this book is arranged
answers family~8. % claim-ok: this page describes how this book is arranged
An application comes where the notion it uses is defined, so the reader meets % claim-ok: this page describes how this book is arranged
the use at the same point as the idea. % claim-ok: this page describes how this book is arranged
An open problem comes where the chapter states the question it belongs to, % claim-ok: this page describes how this book is arranged
and Chapter~\ref{app:questions}, at the front, gathers all of them by family and by class. % claim-ok: this page describes how this book is arranged

A recalled theorem keeps the kind of its source: a proposition stays a % claim-ok: this page describes how this book is arranged
proposition, and a theorem stays a theorem. % claim-ok: this page describes how this book is arranged
The book never states a recalled result again in its own setting. % claim-ok: this page describes how this book is arranged
The specialisation is written after the recall, and it refers to the recall by % claim-ok: this page describes how this book is arranged
number. % claim-ok: this page describes how this book is arranged

\section*{The two appendix chapters}

Appendix~\ref{app:atlas} is the atlas of the literature. % claim-ok: this page describes how this book is arranged
It gives the shape of the bibliography: the works by decade, by language, and by % claim-ok: this page describes how this book is arranged
the kind of derivation-like map they treat. % claim-ok: this page describes how this book is arranged
It is counted again with each edition, so it records what this edition has read % claim-ok: this page describes how this book is arranged
rather than what the field contains. % claim-ok: this page describes how this book is arranged

Chapter~\ref{app:questions}, before Chapter 1, is the collection of question tables. % claim-ok: this page describes how this book is arranged
For each of the eight families it gives, class by class, whether the answer is % claim-ok: this page describes how this book is arranged
known, false with a counterexample, or open, with the source that settles or % claim-ok: this page describes how this book is arranged
asks it. % claim-ok: this page describes how this book is arranged

\section*{Reading paths}

A reader new to the subject can read Part~I in order; it carries the history and % claim-ok: this page describes how this book is arranged
the basic theory, and it fixes the question families. % claim-ok: this page describes how this book is arranged
A reader who works on rings can take Part~II and Part~III, which treat the % claim-ok: this page describes how this book is arranged
associative classes, and consult Part~VI for the analytic ones. % claim-ok: this page describes how this book is arranged
A reader who comes from analysis or operator theory can start with Part~VI and % claim-ok: this page describes how this book is arranged
use Part~II for the algebraic background. % claim-ok: this page describes how this book is arranged
A reader who wants the current frontier can read % claim-ok: this page describes how this book is arranged
Chapter~\ref{app:questions}, which names every open problem with the chapter % claim-ok: this page describes how this book is arranged
that states it, then follow those references back. % claim-ok: this page describes how this book is arranged

\section*{The living edition}

The edition is updated each month. % claim-ok: this page describes how this book is arranged
New work is added where it changes a statement, a table entry or an open % claim-ok: this page describes how this book is arranged
problem. % claim-ok: this page describes how this book is arranged
A work that only confirms what the book already says is not cited for that % claim-ok: this page describes how this book is arranged
alone. % claim-ok: this page describes how this book is arranged
The version number appears on the title page, so that a citation to this book % claim-ok: this page describes how this book is arranged
can name the edition it used. % claim-ok: this page describes how this book is arranged
